\documentclass[11pt]{article}

\usepackage[a4paper,margin=0.9in]{geometry}
\usepackage{amsmath,amssymb,amsthm}
\usepackage{graphicx}
\usepackage[hidelinks]{hyperref}
\usepackage{microtype}

\newtheorem{theorem}{Theorem}
\newtheorem{lemma}{Lemma}
\newtheorem{corollary}{Corollary}
\newcommand{\norm}[1]{\left\lVert #1\right\rVert}
\newcommand{\NA}{\operatorname{NA}}
\newcommand{\N}{\mathbb N}

\title{Property \(\boldsymbol{(\ddagger)}\) of the Dual Forces Reflexivity}
\author{Automatic math-discovery run \texttt{fa\_banach\_001}}
\date{August 11, 2026}

\begin{document}
\maketitle

\section*{Source Problem}

The source is Deepak Gothwal and T.~S.~S.~R.~K. Rao,
\emph{Some Geometric Aspects Related to Lim's Condition},
arXiv:2504.09464, first posted in 2025 and revised in 2026
\cite{GothwalRao}.  On page 11 of the current arXiv PDF, immediately before
Theorem 18, the authors write:

\begin{quote}
``We do not know in general if \(X^*\) has \((\ddagger)\), then \(X\) is
reflexive.''
\end{quote}

They then prove the affirmative answer under the additional assumption that
\(X\) contains no isomorphic copy of \(\ell_1\).

\begin{center}
  \includegraphics[width=\linewidth]{figures/open_problem_crop.png}
\end{center}

\section*{Result Type}

This packet gives a \textbf{candidate full affirmative solution}.  It removes
the source theorem's extra hypothesis \(\ell_1\not\hookrightarrow X\).

\section*{Definition}

A Banach space \(E\) has property \((\ddagger)\) if, whenever
\(0<\varepsilon<2\) and \((f_\alpha)\) is a net of norm-attaining elements of
\(B_{E^*}\) such that
\[
 f_\alpha\xrightarrow{w^*}f,\qquad
 \norm{f_\alpha-f_\beta}\geq\varepsilon\quad(\alpha\ne\beta),
\]
where \(f\in B_{E^*}\) is also norm-attaining, then
\(\norm f\leq 1-\varepsilon/2\).

Here an element of \(E^*\) is norm-attaining if it attains its norm on
\(B_E\).  In the source question \(E=X^*\), so the relevant net lies in
\(B_{X^{**}}\) and the weak-star topology is \(\sigma(X^{**},X^*)\).

\section*{Proof Intuition}

The source paper settles spaces without \(\ell_1\), so the missing case is
when \(\ell_1\) embeds into \(X\).  James's distortion theorem then supplies a
normalized sequence which is a \(0.9\)-copy of the \(\ell_1\) basis.  Its
Stone--\v Cech ultrafilter limits form a weak-star compact family in
\(B_{X^{**}}\), any two of whose points are at distance at least \(1.8\).

Those ultrafilter limits need not themselves be norm-attaining.  This is the
main obstruction, and simply replacing them by Bishop--Phelps approximants
would generally destroy convergence.  The key device is to first translate
the entire separated family so that one chosen limit becomes a nearby
norm-attaining functional.  Each translated point is then perturbed by a
norm-attaining functional with error tending to zero along the net.  The
perturbed net still converges to the chosen norm-attaining limit and retains a
uniform separation.  A common rescaling puts everything in the unit ball,
where property \((\ddagger)\) gives a numerical contradiction.

\section*{Three Lemmas}

\begin{lemma}[Quotient inheritance]\label{lem:quotient}
Let \(E\) have property \((\ddagger)\), and let \(Z\subset E\) be a closed
proximinal subspace.  Then \(E/Z\) has property \((\ddagger)\).
\end{lemma}

\begin{proof}
Identify \((E/Z)^*\) isometrically with \(Z^\perp\subset E^*\).  The weak-star
topology induced by \(E/Z\) agrees on \(Z^\perp\) with the restriction of
\(\sigma(E^*,E)\), because a functional in \(Z^\perp\) depends only on the
coset of its argument.

It remains to check norm attainment.  Suppose \(f\in Z^\perp\) attains its
norm as a functional on \(E/Z\), say at \(x+Z\in S_{E/Z}\).  Proximinality
provides \(z\in Z\) such that
\(\norm{x+z}=\norm{x+Z}=1\).  Since \(f(z)=0\),
\[
  |f(x+z)|=|f(x)|=\norm f.
\]
Thus \(f\) also attains its norm as an element of \(E^*\).  Consequently any
net and limit occurring in the definition for \(E/Z\) also satisfy the
definition's hypotheses in \(E^*\), with the same norms and separation.
Property \((\ddagger)\) for \(E\) gives the required estimate.
\end{proof}

\begin{lemma}[An injective net in \(\N^*\)]\label{lem:net}
There are a point \(p\in\N^*:=\beta\N\setminus\N\) and a net of pairwise
distinct points \((q_d)_{d\in D}\subset\N^*\setminus\{p\}\) such that
\(q_d\to p\).  The directed set may be chosen with a map
\(n:D\to\N\) satisfying \(n(d)\to\infty\).
\end{lemma}

\begin{proof}
Write \(\mathfrak c=2^{\aleph_0}\).  The compact zero-dimensional space
\(\N^*\) has weight at most \(\mathfrak c\), while every nonempty open subset
of \(\N^*\) has cardinality \(2^{\mathfrak c}\).  Indeed, such an open set
contains \(A^*=\overline A^{\,\beta\N}\setminus A\) for some infinite
\(A\subset\N\), and \(A^*\) has cardinality \(2^{\mathfrak c}\).

Fix \(p\in\N^*\), and choose a neighborhood base \(\mathcal B\) at \(p\), of
cardinality at most \(\mathfrak c\), closed under finite intersections.  Put
\(D=\mathcal B\times\N\), ordered by
\[
 (U,n)\preceq(V,m)\quad\Longleftrightarrow\quad V\subset U\text{ and }m\geq n.
\]
Enumerate \(D\) and recursively choose
\(q_{(U,n)}\in U\setminus\{p\}\), never reusing a previously chosen point.
At each stage fewer than or equal to \(\mathfrak c\) points have been used,
whereas \(U\) has cardinality \(2^{\mathfrak c}\), so the recursion succeeds.
The resulting net is injective, converges to \(p\), and its second coordinate
tends to infinity.
\end{proof}

\begin{lemma}[The \(\ell_1\) obstruction]\label{lem:l1}
If a Banach space \(X\) contains an isomorphic copy of \(\ell_1\), then
\(X^*\) does not have property \((\ddagger)\).
\end{lemma}

\begin{proof}
By James's distortion theorem \cite{James,DowlingEtAl}, there is a sequence
\((x_n)\subset S_X\) such that, for every finitely supported scalar sequence
\((a_n)\),
\begin{equation}\label{eq:james}
 \frac9{10}\sum_n|a_n|
 \leq \norm{\sum_n a_nx_n}
 \leq \sum_n|a_n|.
\end{equation}
The map \(n\mapsto x_n\), considered in
\((B_{X^{**}},w^*)\), extends continuously to
\(\Phi:\beta\N\to(B_{X^{**}},w^*)\).  Put \(z_p=\Phi(p)\).

We first record the metric geometry of this compact family.  Given distinct
\(p,q\in\beta\N\), choose \(h=(h_n)\in B_{\ell_\infty}\) whose continuous
extension to \(\beta\N\) has values \(1\) at \(p\) and \(-1\) at \(q\).
On \(\operatorname{span}\{x_n\}\), set
\[
 f_h\!\left(\sum_n a_nx_n\right)
   =\frac9{10}\sum_n a_nh_n.
\]
The lower estimate in \eqref{eq:james} gives \(\norm{f_h}\leq1\), so
Hahn--Banach extends \(f_h\) to an element of \(B_{X^*}\).  Therefore
\begin{equation}\label{eq:separation}
 \norm{z_p-z_q}\geq |(z_p-z_q)(f_h)|=\frac95.
\end{equation}
Taking \(h_n=1\) also gives
\begin{equation}\label{eq:normzp}
 \frac9{10}\leq\norm{z_p}\leq1.
\end{equation}

Fix \(p\in\N^*\).  Apply the Bishop--Phelps theorem to the Banach space
\(X^*\) \cite{BishopPhelps}: choose a norm-attaining
\(w\in X^{**}\) with
\begin{equation}\label{eq:wclose}
 \norm{w-z_p}<\frac1{10}.
\end{equation}
Let \((q_d)_{d\in D}\) and \(n(d)\) be supplied by Lemma~\ref{lem:net}.
Write \(a=w-z_p\).  For every \(d\), Bishop--Phelps again gives a
norm-attaining \(v_d\in X^{**}\) such that
\begin{equation}\label{eq:vd}
 \norm{v_d-(z_{q_d}+a)}<\frac1{10n(d)}.
\end{equation}
Since \(q_d\to p\), we have \(z_{q_d}+a\xrightarrow{w^*}w\); the errors in
\eqref{eq:vd} tend to zero in norm, hence \(v_d\xrightarrow{w^*}w\).

Set
\[
 \phi_d=\frac56v_d,\qquad \phi=\frac56w.
\]
Both the \(\phi_d\)'s and \(\phi\) are norm-attaining.  Moreover,
\eqref{eq:wclose} and \eqref{eq:vd} show that
\[
 \norm{v_d}<1+\frac1{10}+\frac1{10}=\frac65,
\]
so \(\phi_d\in B_{X^{**}}\); also \(\phi\in B_{X^{**}}\), and
\(\phi_d\xrightarrow{w^*}\phi\).  For \(d\ne e\),
\eqref{eq:separation} and
\eqref{eq:vd} yield
\[
 \norm{\phi_d-\phi_e}
 >\frac56\left(\frac95-\frac1{10}-\frac1{10}\right)
 =\frac43.
\]
On the other hand, \eqref{eq:normzp} and \eqref{eq:wclose} give
\[
 \norm\phi=\frac56\norm w
 >\frac56\left(\frac9{10}-\frac1{10}\right)=\frac23.
\]
If \(X^*\) had property \((\ddagger)\), applying it with
\(\varepsilon=4/3\) would instead imply
\(\norm\phi\leq1-(4/3)/2=1/3\), a contradiction.
\end{proof}

\section*{Main Theorem}

\begin{theorem}\label{thm:main}
For every Banach space \(X\), if \(X^*\) has property \((\ddagger)\), then
\(X\) is reflexive.
\end{theorem}

\begin{proof}
Suppose first that \(\ell_1\hookrightarrow X\).  Lemma~\ref{lem:l1} says
that \(X^*\) fails property \((\ddagger)\), contrary to the hypothesis.

It remains to treat \(\ell_1\not\hookrightarrow X\).  We include the source
paper's argument, making the quotient step explicit.  If \(X\) were
nonreflexive, separable determination of reflexivity would give a separable
nonreflexive closed subspace \(Y\subset X\).  Then
\(\ell_1\not\hookrightarrow Y\).  The annihilator \(Y^\perp\) is proximinal
in \(X^*\): if \(x^*\in X^*\), extend \(x^*|_Y\) to
\(g\in X^*\) without increasing its norm; then
\(x^*-g\in Y^\perp\) and
\(\norm g=\operatorname{dist}(x^*,Y^\perp)\).  Since
\(Y^*\cong X^*/Y^\perp\), Lemma~\ref{lem:quotient} shows that \(Y^*\) has
property \((\ddagger)\).

The Odell--Rosenthal theorem \cite{OdellRosenthal} says that every element of
\(B_{Y^{**}}\) is a weak-star sequential limit of a bounded sequence in
\(Y\).  Let \(u\in S_{Y^{**}}\) be norm-attaining.  We may choose a sequence
\((y_n)\subset S_Y\) with \(y_n\xrightarrow{w^*}u\): start with the bounded
sequence from Odell--Rosenthal and normalize it, using a norming functional
for \(u\) to see that its norms tend to one.  If \(u\notin Y\), this sequence
cannot be norm Cauchy, and a subsequence is \(\varepsilon\)-separated for
some \(\varepsilon>0\).  Every \(y_n\), viewed in \(Y^{**}\), is
norm-attaining on \(Y^*\), as is \(u\).  Property \((\ddagger)\) for \(Y^*\)
would give \(1=\norm u\leq1-\varepsilon/2\), impossible.  Hence every
norm-attaining element of \(Y^{**}\) belongs to \(Y\).

By Bishop--Phelps, norm-attaining elements of \(Y^{**}=(Y^*)^*\) are norm
dense in \(Y^{**}\).  Since the canonical copy of \(Y\) is norm closed, it
follows that \(Y=Y^{**}\), contradicting the choice of \(Y\).  Thus \(X\)
is reflexive.
\end{proof}

\begin{corollary}
The open question on page 11 of Gothwal--Rao has an affirmative answer, with
no additional structural hypothesis on \(X\).
\end{corollary}

\section*{Verification Notes}

The proof was audited along the following fault lines.

\begin{itemize}
\item The Stone--\v Cech map is used only into the weak-star compact Hausdorff
ball \(B_{X^{**}}\), so its continuous extension exists by the universal
property of \(\beta\N\).
\item Pairwise separation is obtained from Hahn--Banach extensions whose norm
bound follows in the correct direction from the lower \(\ell_1\) estimate.
\item The net in Lemma~\ref{lem:net} is genuinely injective; a sequence would
not suffice because \(\N^*\) has no nontrivial convergent sequences.
\item Bishop--Phelps is first applied to the desired limit, then to translated
points with errors tending to zero.  This is what preserves both a
norm-attaining limit and weak-star convergence.
\item The common factor \(5/6\) places the whole net in the unit ball.  The
separation lower bound is \(4/3\), while the limit norm is strictly greater
than \(2/3\); property \((\ddagger)\) would force it to be at most \(1/3\).
\item Quotient inheritance is proved directly and does not require a separate
heredity assertion for arbitrary subspaces.
\end{itemize}

The reusable script \texttt{code/check\_constants.py} verifies the numerical
margins, including small perturbations of the three chosen tolerances.  It is
only an arithmetic check; the proof is analytic and topological.

\section*{Deep-Upgrade Audit}

An initial partial route used an asymptotically isometric copy of \(\ell_1\),
for which the ultrafilter points themselves share an exact norming face.
That hypothesis is not automatic: Dowling--Johnson--Lennard--Turett
\cite{DowlingEtAl} construct equivalent norms on \(\ell_1\) containing no
asymptotically isometric copy.  This invalidates the tempting upgrade from
James's theorem.  The translated Bishop--Phelps construction above is the
additional deep attempt that removes the false step and works with the
fixed-error almost-isometric copy actually guaranteed by James.

\section*{Novelty and Literature Check}

The run's registry, solution, attempt, and proof-gap indexes were searched for
arXiv:2504.09464, the exact title, ``Lim's condition'',
``property \((\ddagger)\)'', ``dual'', and ``reflexive''.  No duplicate result
or attempt was found.  The current arXiv source (revised February 2026) still
states the general problem as open and proves only the no-\(\ell_1\) case.

A bounded external search on August 11, 2026 used the exact paper title,
arXiv id, the quoted open sentence, the authors' names, and combinations of
``property \((\ddagger)\)'', ``dual'', and ``reflexive''.  It found the source
record and generic mirrors, but no later paper or independent theorem stated
as resolving this question.  The proof does use classical results of James,
Odell--Rosenthal, and Bishop--Phelps; the particular translation-and-vanishing-
perturbation combination answering the source question was not found.
Accordingly, novelty confidence is moderate rather than definitive.

\section*{Limitations and Human Review Recommendation}

This is a candidate full solution subject to expert review, not a claim of
established publication priority.  The recommended review focus is
Lemma~\ref{lem:net}'s cardinal selection and the two-stage Bishop--Phelps
construction in Lemma~\ref{lem:l1}.  If those steps are accepted, the
constants give a strict contradiction and the source theorem covers the
complementary no-\(\ell_1\) case.

\small
\begin{thebibliography}{9}

\bibitem{GothwalRao}
D. Gothwal and T.~S.~S.~R.~K. Rao,
\emph{Some Geometric Aspects Related to Lim's Condition},
arXiv:2504.09464, 2025; revised 2026.

\bibitem{James}
R.~C. James,
\emph{Uniformly non-square Banach spaces},
Ann. of Math. (2) \textbf{80} (1964), 542--550.

\bibitem{OdellRosenthal}
E. Odell and H.~P. Rosenthal,
\emph{A double-dual characterization of separable Banach spaces containing
\(\ell_1\)},
Israel J. Math. \textbf{20} (1975), 375--384.

\bibitem{BishopPhelps}
E. Bishop and R.~R. Phelps,
\emph{A proof that every Banach space is subreflexive},
Bull. Amer. Math. Soc. \textbf{67} (1961), 97--98.

\bibitem{DowlingEtAl}
P.~N. Dowling, W.~B. Johnson, C.~J. Lennard, and B. Turett,
\emph{The optimality of James's distortion theorems},
Proc. Amer. Math. Soc. \textbf{125} (1997), 167--174.

\end{thebibliography}

\end{document}
